\section{The quantitative theorem}\label{sec:theorem}

Let \(a\in\Q^n\), \(b\in\Q^m\), and let \(\alpha,r>0\) be rational.
Put
\[
 X=\{x:\norm{x-a}\le\alpha\},\qquad
 Y=\{y:\norm{y-b}\le r\},\qquad \Omega=X\times Y.
\]
We first state a normalized theorem for \(H\), and then recover the
usual equation \(F=0\) by setting \(H=AF\).

\begin{definition}[Quantitative \(C^1\) data]\label{def:c1}
A presentation of \(H:\Omega\to\R^m\) consists of primitive recursive
rational evaluators \(\Hhat(z,t)\) and
\(\Jhat(z,t)\) for \(H(z)\) and
\[
 J(z)=D H(z)=[\,U(z)\ V(z)\,],
 \qquad U=D_xH,\quad V=D_yH,
\]
with error at most \(2^{-t}\) in the indicated vector or matrix norm.
It also contains primitive recursive moduli \(\mu,\nu:\N\to\N\)
with the following contracts, including named points in the box:
\begin{align}
 \norm{H(z')-H(z)-J(z)(z'-z)}
  &\le 2^{-p}\norm{z'-z}
 &&\text{if }\norm{z'-z}\le2^{-\mu(p)},\label{eq:taylor}\\
 \norm{J(z')-J(z)}
  &\le 2^{-p}
 &&\text{if }\norm{z'-z}\le2^{-\nu(p)}.\label{eq:duc}
\end{align}
Each named-point contract is shorthand for its finite rational
form. Section~\ref{sec:raw} gives concrete quantifier-free rational
certificates and explicit enlargements of their raw moduli that
imply these contracts. Hence Definition~\ref{def:c1} does not assume
an unformalized real-analysis library.
\end{definition}

\begin{theorem}[Quantitative constructive implicit function theorem]
\label{thm:main}
Suppose the above presentation is given, and a nonnegative rational
\(M\) satisfies, throughout \(\Omega\),
\begin{align}
 \norm{\Id-V(x,y)}&\le\tfrac12,\label{eq:vertical}\\
 \norm{U(x,y)}&\le M,\label{eq:horizontal}\\
 \norm{H(a,b)}&\le r/4,\qquad M\alpha\le r/4.\label{eq:center}
\end{align}
Then there is an explicitly constructed primitive recursive
presentation of a function \(g:X\to Y\) with these properties.
\begin{enumerate}[label=\textup{(\roman*)}]
 \item \(H(x,g(x))=0\), and every \(y\in Y\) satisfying \(H(x,y)=0\)
 equals \(g(x)\).
 \item The following bounds hold:
 \begin{align}
 \norm{g(x)-g(x')}&\le 2M\norm{x-x'},\label{eq:glip}\\
 \norm{y-g(x)}&\le 2\norm{H(x,y)},\label{eq:apost}\\
 \norm{g(x)-b}&\le
 2\bigl(\norm{H(a,b)}+M\norm{x-a}\bigr).\label{eq:location}
 \end{align}
 In particular, if \(H(a,b)=0\), then \(g(a)=b\).
 \item On the interior of \(X\), \(g\) is differentiable and
 \begin{equation}\label{eq:derivative}
 Dg(x)=-V(x,g(x))^{-1}U(x,g(x)).
 \end{equation}
 The same Taylor estimates apply to increments whose endpoints
 remain in \(X\).
 \item Let
 \[
 C=\max(1,2M),\qquad c=\lgup(C),\qquad
 d=\lgup(2+4M).
 \]
 Explicit moduli for the derivative remainder and derivative
 continuity are
 \begin{equation}\label{eq:outputmoduli}
 \mu_g(p)=\mu(p+c+1)+c,\qquad
 \nu_g(p)=\nu(p+d)+c.
 \end{equation}
 More precisely, with both endpoints in \(X\),
 \begin{align*}
 \norm{g(x+h)-g(x)-Dg(x)h}
 &\le2^{-p}\norm h
 &&\text{if }\norm h\le2^{-\mu_g(p)},\\
 \norm{Dg(x')-Dg(x)}
 &\le2^{-p}
 &&\text{if }\norm{x'-x}\le2^{-\nu_g(p)}.
 \end{align*}
\end{enumerate}
For each fixed primitive recursive input presentation, the
finite-observation versions of these conclusions follow in \(\ipra\)
from the corresponding certificate hypotheses. The solution
evaluator, derivative evaluator, and their schedules are primitive
recursive.
\end{theorem}

The theorem allows a certified approximate base zero. For the usual
local graph theorem, one takes \(H(a,b)=0\). In that case
\(\norm{g(x)-b}\le r/2\) on all of \(X\), by
\eqref{eq:center}; thus the graph lies strictly inside the vertical
box.

\begin{corollary}[Original equation]\label{cor:original}
Suppose \(A,B\in\Q^{m\times m}\) satisfy \(AB=BA=\Id\), and
\(H=AF\) has the preceding certificates. Then
\[
 F(x,y)=0\quad\Longleftrightarrow\quad y=g(x)
 \qquad(x\in X,\ y\in Y),
\]
and
\[
 Dg(x)=-D_yF(x,g(x))^{-1}D_xF(x,g(x)).
\]
\end{corollary}

\begin{proof}
\(F=BH\), so the zero predicates are equivalent.
The constructed inverse of \(V=AD_yF\) gives
\((D_yF)^{-1}=V^{-1}A\). Substituting \(U=AD_xF\) in
\eqref{eq:derivative} gives the stated formula. Both identities are
finite matrix calculations with represented entries.
\end{proof}

Requiring an explicit inverse for the rational preconditioner removes
any ambiguity about the passage from \(AF=0\) to \(F=0\).
Section~\ref{sec:conditioning} constructs such a preconditioner from
a positive determinant bound at the base point.
